%% ============================================================
%%  SOLUCIÓN — Ejercicio 1.2: Secciones y formato de texto
%%  Clase 1 — Después de diapositiva "Secciones y jerarquía"
%%  Compilar: F5 (pdflatex)
%% ============================================================

\documentclass[12pt, a4paper]{article}
\usepackage[utf8]{inputenc}
\usepackage[T1]{fontenc}
\usepackage[spanish]{babel}

\title{Salarios y educación en Paraguay}
\author{Tu Nombre}
\date{\today}

\begin{document}
\maketitle

\section{Introducción}

En este trabajo analizamos el efecto de la \textbf{educación} sobre los 
\textit{salarios}. La hipótesis principal es que existe un retorno positivo 
a la escolaridad.

\subsection{Motivación}

Comprender los determinantes del capital humano es \emph{fundamental} para 
el diseño de políticas públicas. En particular, entender cómo cambian los 
retornos a la educación a lo largo del tiempo y entre diferentes grupos 
sociales.

\section*{Notas metodológicas}

Esta sección no tiene número. Usamos datos de panel para controlar por 
\underline{efectos fijos individuales}. Esta metodología nos permite 
eliminar la variación entre individuos y enfocarnos en la variación 
temporal dentro de cada persona.

\end{document}
